\section{From block realizations to the bijective game}
\label{reduction:section}

The general graph construction has two properties with different roles.
Compatible offsets act by bijections on entire block parts, which will
supply Duplicator's announced map. Conversely, a marked zero valuation
receives a legal affine valuation, and simultaneous zero responses agree
on shared readouts. These responses will recover the witness sets of
the scalar-image instance.

We now prove that these two properties suffice. Stating their precise
form here lets the subcubic construction in \Cref{graph:section} use
the same argument after verifying the same conditions. The present
proof uses only the affine blocks of \Cref{offset:section}; it does
not require the later construction.

\begin{proposition}[A realization interface for the game]
\label{prop:game-reduction}
Let a product scalar-image instance have address length \(r\ge3\),
put \(K=r+1\), and form its affine block data. Suppose finite simple
graphs \(W^0,W^1\) of the same positive order have partitions
\[
 V(W^\eps)=\coprod_{B\in\cB}V_B^\eps
 \qquad(\eps\in\{0,1\})
\]
with the following properties.
\begin{enumerate}
\item\label{item:game-local-maps}
For every block \(B\) and offset \(\xi_B\in\cZ_B^1\), there is a
fixed bijection
\(\Psi_{B,\xi_B}:V_B^0\to V_B^1\), depending only on that offset
and the constructed graphs. Whenever offsets on a set of blocks agree
on shared readouts, the union of their maps is an isomorphism of the
induced subgraphs on the corresponding parts.
\item\label{item:game-zero-decoding}
For every \(B\), \(\eps\), and \(z\in\cZ_B^\eps\), there is a
nonempty set \(\mathcal D_B^\eps(z)\subseteq V_B^\eps\); for fixed
\(B,\eps\), these sets are pairwise disjoint. At every history
reached under any winning \(K\)-slot strategy, a marked vertex in
\(\mathcal D_B^0(0)\) is answered by a vertex in
\(\mathcal D_B^1(\zeta_B)\) for some \(\zeta_B\in\cZ_B^1\).
For simultaneous such marks in distinct blocks \(B,C\), the decoded
valuations \(\zeta_B,\zeta_C\) agree on every shared readout.
These assertions apply also when all slots are occupied.
\end{enumerate}
Then
\[
 \text{the instance is consistent}
 \quad\Longleftrightarrow\quad
 \text{Duplicator wins on \(W^0,W^1\) with \(K\) slots}
 \quad\Longleftrightarrow\quad W^0\equiv_r W^1.
\]
\end{proposition}

For the general construction, the query sets are simply
\(\mathcal D_B^\eps(z)=\{(B,z)\}\), and both hypotheses follow from
\Cref{prop:general-realization}. The sets \(\mathcal D_B^\eps(z)\)
allow one valuation to have several query vertices. This flexibility
will let the subcubic construction use a triangle for each valuation.

\begin{proof}
\Cref{lem:wl-game} gives the game/WL equivalence because \(r\ge3\).
We prove the two directions involving consistency from the stated
realization properties.

\paragraph{Witness sets give a winning strategy.}
Suppose the instance has witnesses, and take the nonempty families
$\cH_U$ supplied by \Cref{prop:exact-projection} for every set of blocks
$U$ with $\abs{U}\leq K$.  Recall that their restriction maps satisfy
\eqref{eq:exact-projection} and that every list in such a family agrees on
all shared readouts.

Call a block \emph{active} when its part $V_B^0$ contains a currently
marked first-graph vertex.  Duplicator maintains the following invariant:
if $U$ is the active set, there is a list
\[
 \boldsymbol\xi=(\xi_B)_{B\in U}\in\cH_U
\]
such that every marked vertex $v\in V_B^0$ is paired with
$\Psi_{B,\xi_B}(v)$.  By the local-map hypothesis, the union of these maps is an isomorphism
between the induced subgraphs on the active parts.  In particular, the marked pairs form a partial isomorphism.
Initially $U=\varnothing$, and the invariant holds with the empty list.
After any discards, restrict the list to the remaining active set; exact
projection, in particular its restriction inclusion, keeps it in the
required family.

Suppose a vacant slot is requested.  There are at most $K-1=r$ occupied
slots, so $\abs{U}\leq r$.  For \emph{each} inactive block $B$, choose a
list
\begin{equation}
\label{uniform:preselected-extension}
 \boldsymbol\eta^{\,B}\in\cH_{U\cup\{B\}},
 \qquad
 \boldsymbol\eta^{\,B}|_U=\boldsymbol\xi.
\end{equation}
Such a list exists by the extension direction of exact projection.
All of these choices are made before Spoiler's vertex choice.  They may,
for example, be made using fixed orders on the finite sets of possible
offset lists.  Define a map on the whole first graph by
\[
 f|_{V_B^0}=
 \begin{cases}
  \Psi_{B,\xi_B},&B\in U,\\
  \Psi_{B,\eta^{\,B}_B},&B\notin U.
 \end{cases}
\]
Every piece is a bijection from $V_B^0$ onto $V_B^1$, and these parts
partition the two vertex sets.  Hence $f$ is one bijection of the entire
graphs, which Duplicator announces.

If Spoiler chooses a vertex in an active block, retain the current list.
Otherwise, if the chosen vertex lies in $V_B^0$, use the particular list
$\boldsymbol\eta^{\,B}$ from~\eqref{uniform:preselected-extension} as the
new invariant.  It agrees with every old active-block offset and gives the
announced image of the newly chosen vertex.  No second inactive block
becomes active during this move.  Thus compatibility between the choices
made for two different inactive blocks is unnecessary: only one of those
choices is used to update the invariant.  This proves that Duplicator wins
through every finite history.

\paragraph{A winning strategy gives attained image sets.}
Fix a winning strategy.  For each block \(B\), choose once and for all a
vertex \(z_B\in\mathcal D_B^0(0)\) representing its zero valuation.
This set is nonempty because \(0\in\cZ_B^0\).  A
\emph{zero-query history} is a finite history from empty under the fixed
strategy in which every vertex chosen by Spoiler is one of these $z_B$;
discards are unrestricted.

By the zero-decoding hypothesis, a response to \(z_B\) lies in
\(\mathcal D_B^1(\zeta_B)\) for a uniquely determined legal valuation
$\zeta_B\in\cZ_B^1$.  If zero representatives of distinct blocks are
simultaneously marked, their response valuations agree on every shared
readout. This is the second realization hypothesis, applied at the
current history even when all $K$ slots are occupied. No further test
needs to be executed during a zero-query history. Once a pair is marked,
its actual response vertex and its decoded valuation remain fixed until
that pair is discarded.

For a type $t$ and an allowed full address
$\mathbf a\in\prod_{i=1}^r A_{t,i}$, let
$S_i(t,\mathbf a)$ be the site block with readout
$x_{i,t,a_i,a_{i+1}}$, where indices are cyclic.  These $r$ blocks are
distinct because they have distinct row indices.  We call them the
\emph{site ring} at $(t,\mathbf a)$.  At a zero-query history having
precisely one representative of each of these sites marked, define its
ring sum to be
\begin{equation}
\label{uniform:ring-sum}
 h_t=\sum_{i=1}^r
 \zeta_{S_i(t,\mathbf a)}(x_{i,t,a_i,a_{i+1}})
 \ \in\F^{d_t}.
\end{equation}
There are exactly $r$ occupied slots at such a history, in any arrangement,
and one slot is vacant.  We may write $h_t(\mathfrak h)$ to make the
dependence on the particular history $\mathfrak h$ explicit.

For a real type $P$ and $\mathbf a\in\{0,1\}^r$, let
$\mathfrak H_{P,\mathbf a}$ be the set of all finite zero-query histories
from empty under the fixed strategy that end with precisely the site ring
at $(P,\mathbf a)$ marked, one pair per site in any arrangement of the
slots.  Define
\begin{equation}
\label{uniform:attained-sets}
 Q_{P,\mathbf a}
 =\{h_P(\mathfrak h):\mathfrak h\in\mathfrak H_{P,\mathbf a}\}
 \subseteq\F^{d_P}.
\end{equation}
The definition ranges over \emph{all} such histories under the one fixed
strategy.  Directly querying the $r$ ring sites from empty shows that
every $Q_{P,\mathbf a}$ is nonempty.  We will verify the individual image
equalities for these sets without requiring the strategy to give the same
response on a later visit to a ring.

\paragraph{Transferring one ring across a link.}
Let $\ell=tu$ be a link.  Fix allowed addresses $\mathbf a$ at $t$ and
$\mathbf b$ at $u$ with
$(a_i,b_i)\in\cM_{\ell,i}$ for every $i$.  For this paragraph write
\[
 T_i=S_i(t,\mathbf a),\qquad U_i=S_i(u,\mathbf b),
\]
and let $B_i$ be the square in row $i$ determined by the matches
$(a_i,b_i)$ and $(a_{i+1},b_{i+1})$.  Its horizontal names are the readouts
of $T_i$ and $U_i$, and its vertical names are
\[
 y_i=y_{i,\ell,a_i,b_i},\qquad
 y_{i+1}=y_{i+1,\ell,a_{i+1},b_{i+1}}.
\]
The two endpoint types of a link are distinct, including on a subdivided
loop schema.  In each of these lists the row blocks are distinct, and sites
are distinct from squares.
\Cref{uniform:fig-ring-transfer} displays the two phases and the role of
the spare slot.

Start at any zero-query history with precisely the $T_i$ marked.  Denote
their decoded horizontal vectors by $\alpha_i$, and fix the numerical
value $h_t=\sum_i\alpha_i$ from this starting history.  For
$i=1,\ldots,r$, place $z_{B_i}$ in the vacant slot \emph{before}
discarding the pair at $T_i$.  During that placement there are $r+1=K$
occupied slots.  Since $B_i$ and $T_i$ are then simultaneously marked and
share their $t$-horizontal, the new square's response has that horizontal
equal to $\alpha_i$.

More explicitly, after processing rows $1,\ldots,j$, the marked blocks are
exactly
\[
 B_1,\ldots,B_j,T_{j+1},\ldots,T_r.
\]
For $i\leq j$, the marked square retains $\alpha_i$ on its
$t$-horizontal.  For $i>j$, the original site pair is still present and
retains $\alpha_i$.  This proves by induction that after the first phase
all $r$ squares are marked and their old-side horizontals are the original
vectors $\alpha_i$, not newly selected ring values.

At this all-square history, write $\beta_i$ for the response's
$u$-horizontal in $B_i$.  The adjacent squares $B_{i-1}$ and $B_i$ share
the name $y_i$, so their values on it agree; denote the common value by
$\eta_i$.  The legal valuation equation
\eqref{eq:square-equation}, in version one, reads
\[
 L_{\ell,t}\alpha_i+L_{\ell,u}\beta_i+\eta_i+\eta_{i+1}
   =\delta_i b_\ell.
\]
Summing over all rows cancels each vertical value twice in $\F$.  Since
$\sum_i\delta_i=1$, this gives
\begin{equation}
\label{uniform:transfer-sum}
 L_{\ell,t}h_t+L_{\ell,u}\sum_{i=1}^r\beta_i=b_\ell.
\end{equation}

\begin{figure}[htbp]
\centering
\begin{tikzpicture}[
 x=1cm,y=1cm,line width=.6pt,font=\small,
 stage/.style={draw,rounded corners=2pt,align=center,inner sep=5pt,
               minimum height=1.55cm},
 move/.style={-{Stealth[length=2mm]},shorten >=2pt,shorten <=2pt},
 state/.style={draw,rounded corners=1pt,minimum height=.63cm,
               inner xsep=5pt,align=center}]
 \node[stage,text width=3.2cm] (start) at (1.75,0)
   {\textbf{Starting ring}\\[4pt]
    \(T_1,\ldots,T_r\)\\[3pt]
    \(h_t=\sum_i\alpha_i\)};
 \node[stage,text width=4.1cm] (middle) at (7.65,0)
   {\textbf{All squares marked}\\[4pt]
    \(B_1,\ldots,B_r\)\\[3pt]
    \(\sum_i(\eta_i+\eta_{i+1})=0\)};
 \node[stage,text width=3.2cm] (end) at (13.55,0)
   {\textbf{Ending ring}\\[4pt]
    \(U_1,\ldots,U_r\)\\[3pt]
    \(h_u=\sum_i\beta_i\)};
 \draw[move] (start.east) -- node[above=3pt,align=center,font=\footnotesize]
   {replace\\row by row} (middle.west);
 \draw[move] (middle.east) -- node[above=3pt,align=center,font=\footnotesize]
   {replace\\row by row} (end.west);

 \node[align=center] at (7.65,-1.28)
   {\(L_{\ell,t}\sum_i\alpha_i+
       L_{\ell,u}\sum_i\beta_i=b_\ell\)};

 \node[anchor=west,font=\small\bfseries] at (0,-2.0)
   {First phase: keep \(\alpha_i\) on the \(t\)-horizontal};
 \node[anchor=west,font=\small\bfseries] at (8.0,-2.0)
   {Second phase: keep \(\beta_i\) on the \(u\)-horizontal};
 \node[state] (t) at (.55,-2.85) {\(T_i\)};
 \node[state] (tb) at (3.65,-2.85) {\(T_i,\ B_i\)};
 \node[state] (b) at (6.75,-2.85) {\(B_i\)};
 \draw[move] (t.east) -- node[above=3pt,font=\footnotesize]
   {place \(B_i\)} (tb.west);
 \draw[move] (tb.east) -- node[above=3pt,font=\footnotesize]
   {discard \(T_i\)} (b.west);
 \node[state] (bb) at (8.55,-2.85) {\(B_i\)};
 \node[state] (bu) at (11.65,-2.85) {\(B_i,\ U_i\)};
 \node[state] (u) at (14.75,-2.85) {\(U_i\)};
 \draw[move] (bb.east) -- node[above=3pt,font=\footnotesize]
   {place \(U_i\)} (bu.west);
 \draw[move] (bu.east) -- node[above=3pt,font=\footnotesize]
   {discard \(B_i\)} (u.west);
 \foreach \name/\count in {t/r,tb/{r+1},b/r,bb/r,bu/{r+1},u/r}
   \node[below=4pt,font=\footnotesize] at (\name.south)
     {\(\count\) occupied};
\end{tikzpicture}
\caption{One-link transfer across \(\ell=tu\). Boxes list marked blocks;
arrows are game stages, not output-graph edges. Each row replacement below
keeps the other \(r-1\) marks fixed and places the new zero representative
before discarding the old pair. Simultaneous marks preserve the shared
horizontal value. Each \(B_i\) retains the starting \(t\)-horizontal
\(\alpha_i\). At the common all-square history, \(\beta_i\) is its
\(u\)-horizontal and \(\eta_i\) is the shared value on \(y_i\).
Summing the square equations cancels each \(\eta_i\) twice over \(\F\)
and gives the displayed relation. Each \(U_i\) retains the \(\beta_i\)
from that history. No response must recur after its pair is discarded.}
\label{uniform:fig-ring-transfer}
\end{figure}


Now replace each $B_i$ by $U_i$, again placing the site's zero
representative before removing the square's pair.  The shared
$u$-horizontal copies $\beta_i$ into the new site's response.  After
processing rows $1,\ldots,j$, the marked blocks are
$U_1,\ldots,U_j,B_{j+1},\ldots,B_r$; the new sites retain their copied
$\beta_i$ and the remaining squares retain the values fixed at the
all-square history.  At the end, precisely the $u$-ring is marked and its
sum is $h_u=\sum_i\beta_i$.  Thus the starting and ending sums satisfy
\begin{equation}
\label{uniform:attained-link}
 L_{\ell,t}h_t+L_{\ell,u}h_u=b_\ell.
\end{equation}
Every replacement changes the occupied-slot count from $r$ to $r+1$ and
back to $r$.  The whole transfer uses $2r$ placements and $2r$ discards,
and every chosen vertex is still a fixed zero representative.

This construction also works starting at the $u$-ring: first replace its
sites by the same squares, then replace the squares by the $t$-sites.
The same equation holds with the starting value at $u$ fixed.  No response
is required to recur after its pair has been discarded.

\paragraph{Transferring along a whole comparison wire.}
Fix a schema $e:P\to R$ and an address $\mathbf a\in A_e$.  Along its
wire use $\mathbf a$ at $P$ and at every internal type, and
$\boldsymbol\sigma_e(\mathbf a)$ at $R$.  These are allowed addresses,
and their adjacent coordinates are actual matches at every link: all
but the last link use the restricted identity, while the last uses
$(a_i,\sigma_{e,i}(a_i))$ at coordinate $i$.

Take any $h_P\in Q_{P,\mathbf a}$, choose one witnessing history in
\eqref{uniform:attained-sets}, and continue it by the link transfers just
proved.  The final history is another finite zero-query history, now ending
with exactly the $R,\boldsymbol\sigma_e(\mathbf a)$ ring.  Its sum
$h_R$ therefore belongs to
$Q_{R,\boldsymbol\sigma_e(\mathbf a)}$.  The successive sums satisfy
\eqref{uniform:attained-link}.  Interior link equations equate consecutive
scalar shifts, while the two endpoint equations identify their common
value with the endpoint affine forms.  Consequently,
\[
 \phi_{e,0}(h_P)=\phi_{e,1}(h_R).
\]
This proves
\[
 \phi_{e,0}(Q_{P,\mathbf a})
 \subseteq
 \phi_{e,1}(Q_{R,\boldsymbol\sigma_e(\mathbf a)}).
\]

For the opposite inclusion, start from any history witnessing an element
of $Q_{R,\boldsymbol\sigma_e(\mathbf a)}$ and traverse the wire in reverse.
At its last link, invert each $\sigma_{e,i}$ only on its image: the
specified coordinate $\sigma_{e,i}(a_i)$ has the unique allowed preimage
$a_i$.  All remaining reverse matches are restricted identities.  Thus
the reverse transfer reaches the $P,\mathbf a$ ring and proves the other
inclusion.  No surjectivity of a partial coordinate injection onto the
ambient binary set has been used.

If $P=R$, the endpoint rings are still successive occurrences in the
history.  Their decoded responses may differ, even if the two addresses
are equal.  Both attained values belong to the appropriate set in
\eqref{uniform:attained-sets}, which is all that the image equality requires.
Likewise, separate schemas need no simultaneous choice of endpoint
witnesses.  We have proved \eqref{eq:scalar-images} for every schema and
every address in its domain, establishing consistency.
\end{proof}


\begin{corollary}[Correctness of the general realization]
\label{cor:general-game}
For every product scalar-image instance at address length \(r\ge3\),
its general realization satisfies
\[
 \text{the instance is consistent}\quad\Longleftrightarrow\quad
 G^0\equiv_r G^1.
\]
\end{corollary}
\begin{proof}
Apply \Cref{prop:game-reduction} with the parts, maps and singleton
query sets supplied by \Cref{prop:general-realization}.
\end{proof}
